\documentclass[11pt,a4paper]{article}

\usepackage[utf8]{inputenc}
\usepackage[T1]{fontenc}
\usepackage[italian]{babel}
\usepackage{amsmath}
\usepackage{amssymb}
\usepackage[margin=2.2cm]{geometry}
\usepackage{xcolor}
\usepackage{listings}
\usepackage{enumitem}
\usepackage{booktabs}
\usepackage{longtable}
\usepackage{array}
\usepackage{titlesec}
\usepackage[hidelinks]{hyperref}

\definecolor{codebg}{rgb}{0.96,0.96,0.96}
\definecolor{kw}{rgb}{0.0,0.2,0.6}
\definecolor{cmt}{rgb}{0.3,0.5,0.3}

\lstset{
  basicstyle=\ttfamily\scriptsize,
  backgroundcolor=\color{codebg},
  keywordstyle=\color{kw}\bfseries,
  commentstyle=\color{cmt}\itshape,
  frame=single,
  framesep=4pt,
  breaklines=true,
  showstringspaces=false,
  numbers=left,
  numberstyle=\tiny\color{gray},
  language=C,
  tabsize=4
}

\newcommand{\code}[1]{\texttt{#1}}
\newcommand{\func}[1]{\textbf{\texttt{#1}}}


\titleformat{\section}{\large\bfseries\color{kw}}{\thesection}{0.6em}{}
\titleformat{\subsection}{\normalsize\bfseries}{\thesubsection}{0.6em}{}
\titleformat{\subsubsection}{\normalsize\bfseries\itshape}{\thesubsubsection}{0.6em}{}

\setlist{itemsep=1pt,topsep=2pt}

\title{\textbf{PandOSsh --- Fase 2}\\[2pt]
\large Il Nucleus: inizializzazione, scheduler, eccezioni, interrupt\\
\normalsize Guida dettagliata riga per riga, funzione per funzione}
\author{Materiale di studio per l'interrogazione}
\date{A.A. 2025/2026}

\begin{document}
\maketitle
\tableofcontents
\bigskip
\hrule
\bigskip

% ============================================================
\section{Introduzione generale alla Fase 2}
% ============================================================

La \textbf{Phase 2} costruisce, sopra le strutture dati della Fase 1, il
\textbf{Nucleus}: il kernel vero e proprio di PandOSsh. È il primo livello
del sistema che gira in modo autonomo, gestendo il ciclo di vita completo dei
processi, la loro schedulazione, le eccezioni (incluse le \emph{system
call}) e gli \emph{interrupt}. Il Nucleus gira sempre in kernel-mode, con
gli interrupt disabilitati nei tratti critici, e --- come tutto il progetto
--- senza mai allocare memoria dinamicamente.

Il codice è distribuito su quattro moduli, ciascuno con una responsabilità
precisa:
\begin{center}
\begin{tabular}{@{}llp{7.7cm}@{}}
\toprule
\textbf{File} & \textbf{Ruolo} \\
\midrule
\code{initial.c} & Bootstrap: popola il Pass Up Vector, inizializza la Fase 1
  e le strutture globali, crea il primo processo, avvia lo scheduler. \\
\code{scheduler.c} & Politica di scheduling: round-robin con priorità,
  rilevamento del ``sistema fermo'' (halt/wait/deadlock). \\
\code{exceptions.c} & Punto di ingresso di \emph{tutte} le eccezioni
  sincrone: le 10 system call del Nucleus, i page fault, i program trap, il
  meccanismo di \emph{pass-up-or-die}. \\
\code{interrupts.c} & Gestione degli interrupt asincroni: fine time-slice
  (PLT), tick dello pseudo-clock (Interval Timer), completamento di
  operazioni di I/O sui device. \\
\bottomrule
\end{tabular}
\end{center}

\paragraph{Idea chiave.} Il Nucleus non sa \emph{nulla} della memoria
virtuale o delle U-proc della Fase 3: quando incontra un'eccezione che non
sa gestire (page fault, program trap generico), la ``passa su'' al livello
superiore tramite il meccanismo \code{passUpOrDie}, oppure termina il
processo se questo non ha un livello superiore a cui rivolgersi. Questo
disaccoppiamento è ciò che rende la Fase 3 costruibile \emph{sopra} la Fase
2 \emph{senza modificarla}.

% ============================================================
\section{Vocabolario hardware necessario}
% ============================================================

\begin{description}[style=nextline,leftmargin=0pt]
\item[Kernel vs user mode] Il campo \code{MPP} nel registro \code{status}
  distingue kernel-mode (\code{M}, privilegiato: 0x1800) da user-mode
  (\code{U}: 0). Il Nucleus gira sempre in kernel-mode.
\item[Eccezione vs interrupt] Entrambe deviano il flusso al gestore
  installato nel \emph{Pass Up Vector}. Un'\textbf{eccezione} è sincrona
  (causata dall'istruzione corrente: \code{ecall}, accesso a memoria non
  valido). Un \textbf{interrupt} è asincrono (un device o un timer che
  segnala un evento). Si distinguono guardando il \textbf{bit 31} del
  registro \code{cause} salvato: se acceso, è un interrupt; altrimenti il
  campo basso (\code{ExcCode}) indica il tipo preciso di eccezione.
\item[BIOSDATAPAGE] Indirizzo fisso (\code{0x0FFFF000}) dove il firmware
  salva lo stato completo del processore (\code{state\_t}) nel momento in
  cui avviene un'eccezione/interrupt. Il Nucleus lo interpreta come
  \code{savedState}.
\item[Pass Up Vector] Tabella fissa (\code{0x0FFFF900}) che dice al firmware
  \emph{dove saltare} in caso di eccezione: contiene due coppie
  (indirizzo dell'handler, stack pointer da usare), una per il
  \emph{TLB-Refill event} e una per tutte le altre eccezioni.
\item[TOD / PLT / Interval Timer] \code{TOD} (Time Of Day) è l'orologio
  assoluto, letto con \code{STCK}, usato per misurare il tempo CPU. Il
  \textbf{PLT} (Processor Local Timer) è il timer che genera l'interrupt di
  fine time-slice (preemption), impostato con \code{setTIMER}.
  L'\textbf{Interval Timer} è il timer di sistema che batte lo
  \emph{pseudo-clock} ogni \code{PSECOND} (100\,ms), riarmato con
  \code{LDIT}.
\end{description}

\subsection{Le primitive della libreria \code{uriscv/liburiscv.h}}
Tutte le funzioni in \code{MAIUSCOLO} (\code{STCK}, \code{LDST}, \ldots)
che compaiono nel codice dei tre documenti \textbf{non sono funzioni del
progetto}: sono primitive fornite dalla SDK dell'emulatore
(\code{uriscv/liburiscv.h}), implementate in assembly, che eseguono
un'istruzione macchina privilegiata o accedono direttamente ai registri
del processore. Vale la pena avere un riferimento unico di cosa fa
ciascuna, dato che vengono usate in continuazione senza ulteriori
spiegazioni nei listati:

\begin{center}
\begin{longtable}{@{}llp{7.8cm}@{}}
\toprule
\textbf{Primitiva} & \textbf{Categoria} & \textbf{Cosa fa} \\
\midrule
\code{SYSCALL(n,a1,a2,a3)} & syscall & Esegue l'istruzione \code{ecall}
  con \code{n} in \code{a0} e gli argomenti in \code{a1--a3}: genera
  l'eccezione che porta a \code{syscallHandler} (Nucleus) o
  \code{supSyscallHandler} (Support Level). Ritorna il valore lasciato in
  \code{a0} dal gestore. \`E l'\emph{unico} modo con cui un processo
  comunica con un livello sottostante. \\
\code{STCK(v)} & tempo & Legge il TOD (Time-Of-Day) corrente nella
  variabile \code{v}: la base di ogni misura di tempo (contabilità CPU,
  \code{GETTIME}, \code{CLOCKWAIT}). \\
\code{STST(statep)} & stato & \emph{Store State}: salva lo stato
  \emph{corrente} del processore in \code{*statep}. Usata solo dal test
  driver \code{p2test.c} per catturare un template di stato da cui partire
  nel costruire gli stati iniziali dei processi di test. \\
\code{LDST(statep)} & stato & \emph{Load State}: carica \code{*statep} nel
  processore e \textbf{riprende l'esecuzione da lì} --- non ritorna mai al
  chiamante. \`E l'istruzione con cui lo scheduler avvia/riprende un
  processo, e con cui ogni gestore di eccezione ``chiude'' la gestione
  restituendo il controllo. \\
\code{LDCXT(sp,status,pc)} & pass-up & \emph{Load Context}: come
  \code{LDST}, ma invece di un intero \code{state\_t} prende separatamente
  i tre valori di un \code{context\_t} (si veda il documento di Fase 1) e
  salta lì --- è il meccanismo con cui \code{passUpOrDie} trasferisce il
  controllo al Support Level. \\
\code{WAIT()} & CPU & Sospende la CPU a basso consumo finché non arriva un
  interrupt (usata dallo scheduler quando non c'è nulla di pronto ma
  qualcuno è soft-blocked). \\
\code{PANIC()} & arresto & Arresta la macchina virtuale segnalando un
  errore fatale/irrecuperabile (deadlock, invariante violato). \\
\code{HALT()} & arresto & Arresta la macchina virtuale in modo
  ``pulito'', quando \code{processCount == 0}: è la conclusione normale
  del sistema. \\
\code{setTIMER(v)} & timer & Carica il valore \code{v} nel PLT (Processor
  Local Timer): allo scadere genera l'interrupt \code{excCode == 7}.
  Caricato con \code{NEVER} per disarmarlo. \\
\code{LDIT(v)} & timer & Carica il valore \code{v} nell'Interval Timer
  (riarmato a ogni tick con \code{PSECOND}, cioè ogni 100\,ms). \\
\code{getSTATUS()} / \code{setSTATUS(v)} & registri & Leggono/scrivono il
  registro \code{status} \emph{corrente} del processore (diverso da
  \code{state\_t->status}, che è una \emph{copia salvata}): usate per
  abilitare/disabilitare bit specifici a runtime (es. il PLT prima di un
  \code{WAIT}, o gli interrupt attorno a un aggiornamento di TLB in Fase
  3). \\
\code{setMIE(v)} & registri & Scrive la maschera delle sorgenti di
  interrupt abilitate (es. \code{MIE\_ALL \& \textasciitilde
  MIE\_MTIE\_MASK} per abilitare tutto tranne il PLT). \\
\code{TLBCLR()} & TLB & Invalida (azzera) tutte le entry del TLB. \\
\code{TLBWR()} & TLB & \emph{TLB Write Random}: scrive una nuova
  traduzione (da \code{ENTRYHI}/\code{ENTRYLO}) in una entry del TLB
  scelta automaticamente dall'hardware. Usata sia dal TLB-Refill handler
  sia dal Pager (Fase 3) per installare una traduzione. \\
\code{setENTRYHI(v)} / \code{setENTRYLO(v)} & TLB & Impostano i registri
  di appoggio \code{EntryHI}/\code{EntryLO} \emph{prima} di una
  \code{TLBWR}: rispettivamente VPN+ASID e PFN+bit di controllo
  (\code{V}, \code{D}) della traduzione da scrivere. \\
\bottomrule
\end{longtable}
\end{center}
Nessuna di queste primitive è definita nel codice del progetto: sono
dichiarate \code{extern} in \code{uriscv/liburiscv.h} e la loro
implementazione reale (poche istruzioni assembly ciascuna) vive
nell'installazione della SDK dell'emulatore, non nei sorgenti consegnati.

% ============================================================
\section{\code{phase2/headers/globals.h} --- le variabili globali del Nucleus}
% ============================================================

\begin{lstlisting}
extern int processCount;
extern int softBlockCount;
extern struct list_head readyQueue;
extern pcb_t *currentProcess;
extern pcb_t *yieldedProcess;
extern pcb_t *activeProcs[MAXPROC];

#define DEV_SEM_BASE(line, dev) (((line) - 3) * 8 + (dev))
#define TERM_TX_SEM(dev)  (32 + (dev))
#define TERM_RX_SEM(dev)  (40 + (dev))
#define PSEUDOCLK_SEM           48
#define TOT_SEMS                49

extern int devSems[TOT_SEMS];
extern cpu_t startTOD;
\end{lstlisting}

Tutte queste variabili sono \emph{definite} (con lo storage reale) in
\code{initial.c} e dichiarate \code{extern} qui per essere condivise dagli
altri tre moduli.

\begin{itemize}
\item \code{processCount}: numero di processi \textbf{vivi} (creati e non
  ancora terminati), in qualunque stato (pronto, in esecuzione, bloccato).
  Quando arriva a 0 significa che tutto il lavoro del sistema è concluso.
\item \code{softBlockCount}: numero di processi bloccati in attesa di un
  evento \emph{esterno asincrono} (un device o lo pseudo-clock). Serve allo
  scheduler per distinguere ``sto aspettando un interrupt'' da ``sono in
  deadlock''.
\item \code{readyQueue}: la coda dei processi pronti, ordinata per priorità
  (vedi \code{insertProcQ} in Fase 1).
\item \code{currentProcess}: puntatore al PCB del processo in esecuzione in
  questo momento, oppure \code{NULL} se nessun processo sta girando (il
  Nucleus stesso è ``tra'' due processi).
\item \code{yieldedProcess}: puntatore temporaneo usato per implementare la
  syscall \code{YIELD} (SYS10): tiene il processo che ha appena ceduto
  volontariamente la CPU, in modo che lo scheduler possa decidere se farlo
  ripartire subito o dargli priorità inferiore rispetto a un altro processo
  pronto di pari o maggiore priorità.
\item \code{activeProcs[MAXPROC]}: array di puntatori a \emph{tutti} i
  processi vivi (non solo quelli pronti). \`E una struttura \textbf{aggiunta
  rispetto alla specifica minima} di Pandos, introdotta per due motivi
  precisi: (1) la syscall \code{TERMPROCESS} deve poter trovare un processo
  \emph{per PID} anche se questo è bloccato (quindi non nella
  \code{readyQueue}); (2) lo scheduler deve poter distinguere in modo
  affidabile la condizione di deadlock. Il costo è $O(\text{MAXPROC})$ per
  ricerca/inserimento/rimozione, accettabile dato il limite fisso di 20
  processi.
\item \code{devSems[49]}: array dei semafori dei device. I primi 48 elementi
  coprono le 5 linee di interrupt esterne (disk, flash, ethernet, printer,
  terminal) $\times$ 8 device ciascuna secondo la formula
  \code{DEV\_SEM\_BASE(line, dev) = (line-3)*8 + dev}; per i terminali,
  che hanno due sotto-canali indipendenti (trasmissione e ricezione), si
  usano indici dedicati \code{TERM\_TX\_SEM}/\code{TERM\_RX\_SEM} (che di
  fatto sovrascrivono/estendono lo spazio della linea terminale, indici
  32--47). L'ultimo elemento, indice 48 (\code{PSEUDOCLK\_SEM}), è il
  semaforo dello pseudo-clock.
\item \code{startTOD}: valore del TOD registrato all'ultimo punto in cui si
  è iniziato a ``contare'' il tempo CPU del processo corrente (dispatch,
  syscall, interrupt): confrontato con il TOD attuale permette di calcolare
  quanti cicli attribuire a \code{p\_time}.
\end{itemize}

% ============================================================
\section{\code{phase2/debug.h} --- macro di debug condizionale}
% ============================================================

Prima di entrare nei quattro moduli veri e propri, vale la pena conoscere
\code{debug.h}: è incluso da \emph{tutti} (\code{initial.c},
\code{scheduler.c}, \code{exceptions.c}, \code{interrupts.c}) e definisce
il meccanismo di stampa usato per tracciare l'esecuzione del Nucleus a
runtime.

\begin{lstlisting}
#define DEBUG_INIT 0
#define DEBUG_SCHED 0
#define DEBUG_EXC 0
#define DEBUG_INT 0

#define DEBUG_ENABLED 1

#if DEBUG_ENABLED
#define DELAY() for (volatile int i = 0; i < 60; i++)
static inline void debug_print(const char *msg) {
    unsigned int *command = (unsigned int *)(0x10000254 + 3*4);
    while (*msg != '\0') {
        *command = 2 | (((unsigned int)*msg) << 8);
        DELAY();
        msg++;
    }
}
static inline void debug_hex(const char *label, unsigned int val) {
    /* stampa label seguita dal valore in esadecimale a 8 cifre */
}
#else
#define debug_print(msg) ((void)0)
#define debug_hex(label, val) ((void)0)
#endif
#endif
\end{lstlisting}

\begin{itemize}
\item \textbf{Due livelli di controllo indipendenti}: un interruttore
  \emph{globale} (\code{DEBUG\_ENABLED}) e quattro interruttori
  \emph{per modulo} (\code{DEBUG\_INIT}, \code{DEBUG\_SCHED},
  \code{DEBUG\_EXC}, \code{DEBUG\_INT}). Ogni modulo, in testa al proprio
  file, ridefinisce le macro \code{IDBG}/\code{SDBG}/\code{EDBG} (con nome
  diverso solo per evitare collisioni tra moduli) in base al proprio flag:
  se attivo, espandono a \code{debug\_print}/\code{debug\_hex}; se
  disattivo, espandono a \code{((void)0)} (un'istruzione nulla, eliminata
  dal compilatore, così da non lasciare \emph{alcun} costo a runtime quando
  il debug è spento). Nella versione consegnata tutti e quattro i flag di
  modulo sono a \code{0}: il debug è disattivabile modulo per modulo senza
  ricompilare tutto il progetto, ma di default è silenzioso.
\item \code{debug\_print(msg)}: scrive una stringa \textbf{direttamente sui
  registri hardware del terminale 0} (indirizzo \code{0x10000254 + 3*4},
  cioè il registro di comando della trasmissione), un carattere alla
  volta, con un comando \code{PRINTCHR} (\code{2}) composto a mano
  (carattere nel byte alto, opcode nel byte basso) --- \textbf{non} passa
  mai da \code{DOIO}: è pensato per stampare diagnostica \emph{anche prima}
  che il Nucleus sia in grado di gestire syscall e interrupt correttamente
  (utile per debuggare proprio \code{initial.c} durante il boot), o
  comunque senza dipendere dal meccanismo che si starebbe cercando di
  diagnosticare.
\item \code{DELAY()}: un ciclo di attesa attiva (\emph{busy-wait}) da 60
  iterazioni tra un carattere e l'altro, necessario perché la scrittura sul
  registro di comando non è istantanea: senza una pausa minima, caratteri
  scritti troppo rapidamente in successione andrebbero persi (nessun
  meccanismo di sincronizzazione reale è coinvolto, è un semplice ritardo
  empirico).
\item \code{debug\_hex(label, val)}: come \code{debug\_print}, ma stampa
  anche un valore a 32 bit convertito manualmente in 8 cifre esadecimali
  (utile per stampare indirizzi di semafori, PID, valori di registri senza
  dover comporre a mano la stringa a ogni chiamata).
\end{itemize}
\`E lo strumento di debug \emph{effettivamente usato} durante lo sviluppo
del Nucleus (a differenza di \code{klog.c}, discusso nel documento di Fase
1, che è compilato ma mai invocato): le chiamate a \code{IDBG\_HEX},
\code{SDBG}, \code{EDBG\_HEX}, ecc.\ visibili nei listati dei moduli
successivi sono per l'appunto invocazioni di queste macro.

% ============================================================
\section{\code{phase2/initial.c} --- inizializzazione del sistema}
% ============================================================

\subsection{Variabili globali definite qui}
\begin{lstlisting}
int              processCount;
int              softBlockCount;
struct list_head readyQueue;
pcb_t           *currentProcess;
int              devSems[TOT_SEMS];
cpu_t            startTOD;
pcb_t           *activeProcs[MAXPROC];
pcb_t           *yieldedProcess = NULL;
\end{lstlisting}
Questo è l'\emph{unico} punto del progetto in cui queste variabili vengono
effettivamente allocate (storage reale); ovunque altrove sono \code{extern}.

\subsection{\func{main} --- la funzione di bootstrap}
\begin{lstlisting}
int main(void) {
    passupvector_t *passUpVec = (passupvector_t *) PASSUPVECTOR;
    memaddr ramtop;
    RAMTOP(ramtop);

    passUpVec->tlb_refill_handler  = (memaddr) uTLB_RefillHandler;
    passUpVec->tlb_refill_stackPtr = ramtop;
    passUpVec->exception_handler   = (memaddr) exceptionHandler;
    passUpVec->exception_stackPtr  = ramtop - PAGESIZE;

    initPcbs();
    initASL();

    processCount   = 0;
    softBlockCount = 0;
    currentProcess = NULL;
    mkEmptyProcQ(&readyQueue);

    for (int i = 0; i < TOT_SEMS; i++) devSems[i] = 0;

    STCK(startTOD);

    for (int i = 0; i < MAXPROC; i++)
        activeProcs[i] = NULL;

    LDIT(PSECOND);

    pcb_t *testPcb = allocPcb();
    if (testPcb == NULL) PANIC();

    testPcb->p_s.status      = MSTATUS_MIE_MASK | MSTATUS_MPIE_MASK | MSTATUS_MPP_M;
    testPcb->p_s.mie         = MIE_ALL;
    testPcb->p_s.reg_sp      = ramtop - (2 * PAGESIZE);
    testPcb->p_s.pc_epc      = (memaddr) test;
    testPcb->p_prio          = PROCESS_PRIO_LOW;

    activeProcs[0] = testPcb;
    insertProcQ(&readyQueue, testPcb);
    processCount = 1;

    scheduler();
    return 0;
}
\end{lstlisting}

Passo per passo:

\begin{enumerate}
\item \textbf{Calcolo di \code{ramtop}.} La macro \code{RAMTOP(ramtop)}
  legge dai registri di configurazione della macchina (\code{RAMBASEADDR},
  \code{RAMBASESIZE}) l'indirizzo del \emph{primo byte oltre la fine della
  RAM} disponibile.

\item \textbf{Popolamento del Pass Up Vector.} Il Nucleus registra se stesso
  come gestore \emph{unico} di tutte le eccezioni hardware:
  \begin{itemize}
  \item \code{exception\_handler = exceptionHandler}: qualunque eccezione
    ``normale'' (syscall, program trap, page fault non-refill) salta a
    \code{exceptionHandler} (in \code{exceptions.c});
  \item \code{tlb\_refill\_handler = uTLB\_RefillHandler}: il caso speciale
    di TLB-Refill (percorso veloce, vedi Fase 3) salta a un handler
    dedicato;
  \item \textbf{scelta implementativa importante}: gli stack di questi due
    gestori sono impostati sui \emph{due frame più alti della RAM}
    (\code{ramtop} per il TLB-Refill, \code{ramtop - PAGESIZE} per le
    eccezioni), \emph{non} su un unico indirizzo fisso condiviso
    (\code{KERNELSTACK}). Motivo: se si usasse un solo stack, un'eccezione
    che avviene \emph{durante} la gestione di un'altra (eccezioni annidate,
    possibili perché nulla vieta a un TLB-Refill di avvenire mentre gira
    l'exception handler) sovrascriverebbe lo stack dell'altro handler,
    corrompendone lo stato. Isolarli su due frame fisici separati elimina
    questo rischio.
  \end{itemize}

\item \textbf{Inizializzazione della Fase 1.} \code{initPcbs()} e
  \code{initASL()} preparano le free list dei PCB e dei descrittori di
  semaforo (vedi documento di Fase 1).

\item \textbf{Azzeramento delle strutture globali del Nucleus}:
  \code{processCount}, \code{softBlockCount}, \code{currentProcess = NULL},
  \code{readyQueue} vuota, tutti i 49 \code{devSems} a 0 (semantica a
  contatore: valore iniziale 0, non negativo, perché nessun processo è
  ancora bloccato su di essi).

\item \textbf{Lettura del TOD reale} (\code{STCK(startTOD)}) invece di
  azzerare \code{startTOD}. Se si azzerasse, il primo processo creato
  ``erediterebbe'' come tempo CPU l'intero intervallo trascorso dall'avvio
  della macchina virtuale fino a questo punto del bootstrap (che può essere
  non trascurabile), falsando la contabilità di \code{GETTIME}.

\item \code{activeProcs[i] = NULL} per tutti gli slot: nessun processo vivo
  ancora.

\item \textbf{Avvio dell'Interval Timer}: \code{LDIT(PSECOND)} programma il
  primo tick dello pseudo-clock tra 100\,000 microsecondi (100\,ms,
  \code{PSECOND} in \code{const.h}).

\item \textbf{Creazione del processo iniziale} (`test`, il punto di ingresso
  della Fase 2/3 --- in Fase 2 puro esegue i test P2, in Fase 3 è
  l'\code{InstantiatorProcess}): si alloca un PCB con \code{allocPcb()}
  (fallimento $\to$ \code{PANIC()}, situazione impossibile dato che la free
  list è appena stata popolata con 20 slot), e si configura manualmente lo
  stato iniziale del processore:
  \begin{itemize}
  \item \code{status}: interrupt abilitati (\code{MIE}), interrupt
    ``precedenti'' abilitati (\code{MPIE}, ripristinati automaticamente al
    primo \code{LDST}), \code{MPP} = kernel-mode (\code{MPP\_M}) --- il
    processo test gira in kernel-mode, non è una U-proc;
  \item \code{mie = MIE\_ALL}: tutte le sorgenti di interrupt abilitate a
    livello di maschera;
  \item \code{reg\_sp = ramtop - 2*PAGESIZE}: lo stack del processo test è
    posizionato \emph{sotto} i due frame riservati agli handler del kernel
    (che occupano \code{ramtop} e \code{ramtop-PAGESIZE}), evitando ogni
    sovrapposizione;
  \item \code{pc\_epc = (memaddr) test}: il processo parte eseguendo la
    funzione \code{test} (definita in \code{p2test.c} per la Fase 2 pura, o
    in \code{phase3/initProc.c} per la Fase 3, a seconda del target di
    compilazione);
  \item \code{p\_prio = PROCESS\_PRIO\_LOW}: priorità bassa (0).
  \end{itemize}
  Il PCB viene poi registrato in \code{activeProcs[0]}, inserito nella
  \code{readyQueue}, e \code{processCount} viene impostato a 1.

\item \textbf{Avvio dello scheduler}: \code{scheduler()} non ritorna mai
  (in condizioni normali termina l'esecuzione con \code{HALT()}); la riga
  \code{return 0;} successiva è dead code raggiunto solo in caso di bug
  grave.
\end{enumerate}

Da notare la scelta di \emph{non} duplicare in \code{main} la logica di
selezione ``chi eseguire per primo'': si crea un solo processo pronto e si
lascia allo \code{scheduler()} il compito di estrarlo e avviarlo, mantenendo
tutta la logica di dispatch in un unico posto.

% ============================================================
\section{\code{phase2/scheduler.c} --- lo scheduler}
% ============================================================

\begin{lstlisting}
void scheduler(void) {

    if (yieldedProcess != NULL) {
        pcb_t *y = yieldedProcess;
        yieldedProcess = NULL;
        if (!emptyProcQ(&readyQueue) &&
            headProcQ(&readyQueue)->p_prio >= y->p_prio) {
            insertProcQ(&readyQueue, y);
        } else {
            currentProcess = y;
            STCK(startTOD);
            setTIMER(TIMESLICE * (*((cpu_t *) TIMESCALEADDR)));
            LDST(&currentProcess->p_s);
        }
    }

    if (!emptyProcQ(&readyQueue)) {
        currentProcess = removeProcQ(&readyQueue);
        STCK(startTOD);
        setTIMER(TIMESLICE * (*((cpu_t *) TIMESCALEADDR)));
        LDST(&currentProcess->p_s);
    }

    currentProcess = NULL;

    if (processCount == 0) {
        HALT();
    }

    if (softBlockCount > 0) {
        setMIE(MIE_ALL & ~MIE_MTIE_MASK);
        unsigned int status = getSTATUS();
        status |= MSTATUS_MIE_MASK;
        setSTATUS(status);
        WAIT();
    }

    PANIC();
}
\end{lstlisting}

La politica implementata è un \textbf{round-robin con priorità}: tra i
processi pronti si sceglie sempre il più prioritario (grazie
all'ordinamento già garantito da \code{insertProcQ} in Fase 1); a parità di
priorità, l'ordine di arrivo (FIFO) decide chi va per primo, e ciascuno gira
per al massimo una \emph{time slice} prima di essere rimesso in coda (dal
gestore dell'interrupt PLT).

\subsection{Passo 0 --- gestione dello YIELD}
\label{sub:yield-scheduler}
Se \code{yieldedProcess} non è \code{NULL} (un processo ha appena eseguito
la syscall \code{YIELD}, SYS10), lo scheduler decide se rimetterlo
\emph{immediatamente} in coda oppure farlo ripartire subito:
\begin{itemize}
\item se la ready queue \emph{non} è vuota e il processo in testa ha
  priorità $\ge$ a quella dello yielder $\to$ lo yielder \emph{cede
  davvero}: viene reinserito in coda con \code{insertProcQ} (che lo
  posizionerà secondo la sua priorità), e si prosegue al passo successivo
  (che estrarrà dalla coda il processo più prioritario, possibilmente lo
  stesso yielder se nel frattempo nessuno lo ha superato --- ma tipicamente
  un altro);
\item altrimenti (coda vuota, oppure lo yielder è comunque il più
  prioritario disponibile) \code{currentProcess} diventa di nuovo lo
  yielder: si registra il TOD, si riarma il PLT con una nuova time slice
  completa, e si riparte con \code{LDST}. \emph{Senza} questo ramo, uno
  \code{YIELD} in un sistema con un solo processo pronto risulterebbe in un
  \code{removeProcQ} seguito da un nuovo giro immediato comunque, ma la
  gestione esplicita rende il comportamento uniforme ed evita una doppia
  ri-scansione della coda.
\end{itemize}

\subsection{Passo 1 --- dispatch di un processo pronto}
Se la ready queue non è vuota, se ne estrae la testa
(\code{removeProcQ}, $O(1)$, rispetta la priorità), si registra
\code{startTOD} (per la contabilità del tempo CPU che seguirà), si carica
nel PLT una nuova \textbf{time slice} pari a \code{TIMESLICE *
TIMESCALEADDR} (\code{TIMESLICE = 5\,000\,000}, moltiplicato per il fattore
di scala del TOD letto dal registro hardware \code{TIMESCALEADDR}, per
ottenere il valore corretto nell'unità di tempo della macchina), e si
riprende l'esecuzione del processo con \code{LDST} (Load State: ripristina
tutti i registri e il PC dal \code{p\_s} salvato, e non ritorna mai al
chiamante --- è un vero e proprio salto).

\subsection{Passi 2--4 --- il sistema è ``fermo''}
Se si arriva qui, non c'era nessun processo pronto da eseguire
(\code{currentProcess} viene esplicitamente riportato a \code{NULL}): lo
scheduler distingue tre possibili situazioni, in ordine:
\begin{enumerate}
\item \textbf{\code{processCount == 0}}: non esiste più nessun processo
  vivo, tutto il lavoro del sistema è concluso $\to$ \code{HALT()} ferma
  la macchina virtuale.
\item \textbf{\code{softBlockCount > 0}}: esistono processi vivi ma
  \emph{tutti} sono in attesa di un evento esterno (I/O o pseudo-clock).
  Il Nucleus non ha nulla da eseguire ma \emph{non} deve fermarsi: entra in
  \code{WAIT()} (istruzione che sospende la CPU a basso consumo finché non
  arriva un interrupt). Prima di farlo, abilita tutti gli interrupt
  \textbf{tranne} il PLT (\code{MIE\_ALL \& \textasciitilde
  MIE\_MTIE\_MASK}): non ha senso ricevere un interrupt di fine time-slice
  quando nessun processo sta effettivamente ``consumando'' uno slice, e si
  evita così un interrupt PLT spurio (che complicherebbe inutilmente la
  gestione, dato che \code{currentProcess} sarebbe \code{NULL}).
\item \textbf{Altrimenti} (processi vivi, nessuno pronto, nessuno
  soft-blocked): è un vero \textbf{deadlock} --- non esiste alcun evento
  futuro (né un dispatch, né un interrupt atteso) che possa far progredire
  il sistema. Si chiama \code{PANIC()}, che arresta la macchina segnalando
  l'errore.
\end{enumerate}

% ============================================================
\section{\code{phase2/exceptions.c} --- eccezioni e system call}
% ============================================================

Questo è il modulo più corposo del Nucleus. Gestisce: il TLB-Refill event
(percorso condizionale per la Fase 3), il dispatch di tutte le eccezioni
sincrone, l'implementazione delle 10 syscall del Nucleus, la terminazione
ricorsiva dei processi, e il meccanismo di pass-up-or-die.

\subsection{Funzioni di utilità (\code{static})}

\begin{lstlisting}
static void copyState(state_t *dst, state_t *src) {
    unsigned int *d = (unsigned int *) dst;
    unsigned int *s = (unsigned int *) src;
    for (int i = 0; i < STATE_T_SIZE_IN_BYTES / WORDLEN; i++)
        d[i] = s[i];
}
\end{lstlisting}
\func{copyState}: copia byte-a-byte (in realtà word-a-word, 4 byte alla
volta) l'intero contenuto di uno \code{state\_t} da \code{src} a
\code{dst}, trattando le due strutture come array di interi di dimensione
\code{STATE\_T\_SIZE\_IN\_BYTES / WORDLEN} (148/4 = 37 word). È usata al
posto di un'assegnazione diretta \code{*dst = *src} probabilmente per
compatibilità/chiarezza esplicita a livello di word; il risultato è
equivalente. Serve ogni volta che occorre salvare lo stato del processore
letto da \code{BIOSDATAPAGE} dentro il PCB del processo corrente
(\code{currentProcess->p\_s}), tipicamente prima di bloccare o rischedulare
il processo.

\begin{lstlisting}
static void updateCPUTime(void) {
    if (currentProcess) {
        cpu_t now;
        STCK(now);
        currentProcess->p_time += now - startTOD;
        startTOD = now;
    }
}
\end{lstlisting}
\func{updateCPUTime}: implementa la \textbf{contabilità del tempo CPU}. A
ogni ``punto critico'' (ingresso in una syscall, in un interrupt, in un
blocco) si legge il TOD corrente (\code{STCK}), si calcola il delta rispetto
all'ultimo \code{startTOD} registrato, lo si somma a \code{p\_time} del
processo corrente, e si aggiorna \code{startTOD} al nuovo valore. Questo
schema, applicato sistematicamente in \emph{ogni} punto in cui il controllo
lascia il processo corrente (anche solo temporaneamente, come durante la
gestione di una syscall), garantisce che nessun ciclo di CPU vada perso o
venga contato due volte, anche in presenza di eccezioni annidate o di più
syscall consecutive nello stesso quanto di tempo.

\begin{lstlisting}
static int isDeviceSemaphore(int *semAdd) {
    int *base = &devSems[0];
    int *top  = &devSems[TOT_SEMS];
    if (semAdd < base || semAdd >= top) return 0;
    int idx = (int)(semAdd - base);
    if (idx == PSEUDOCLK_SEM) return 0;
    return 1;
}
\end{lstlisting}
\func{isDeviceSemaphore}: distingue, dato l'indirizzo di un semaforo, se è
un semaforo di \emph{device} (vero) oppure un semaforo ``normale'' (di
sincronizzazione tra processi, oppure lo pseudo-clock: falso). Verifica
prima che l'indirizzo cada dentro l'array \code{devSems} (aritmetica dei
puntatori: confronto tra indirizzi), poi esclude esplicitamente l'ultimo
indice (\code{PSEUDOCLK\_SEM}), che pur essendo fisicamente dentro
\code{devSems} rappresenta lo pseudo-clock, contato separatamente. Questa
distinzione è cruciale per \code{blockCurrentProcess}: il contatore
\code{softBlockCount} deve incrementare per i device \emph{ma anche} per lo
pseudo-clock (che infatti viene gestito a parte, incrementando
esplicitamente in \code{CLOCKWAIT}) --- la funzione qui serve invece a
\emph{escludere} lo pseudo-clock dal doppio conteggio quando ci si blocca
tramite il percorso generico \code{blockCurrentProcess}.

\begin{lstlisting}
static void blockCurrentProcess(int *sem) {
    if (!currentProcess) PANIC();
    cpu_t now;
    STCK(now);
    currentProcess->p_time += now - startTOD;
    startTOD = now;
    currentProcess->p_semAdd = sem;
    if (isDeviceSemaphore(sem)) {
        softBlockCount++;
    }
    insertBlocked(sem, currentProcess);
    currentProcess = NULL;
    scheduler();
}
\end{lstlisting}
\func{blockCurrentProcess}: centralizza la logica di blocco di
\code{currentProcess} su un semaforo qualsiasi. Aggiorna il tempo CPU
(stessa logica di \code{updateCPUTime}, duplicata inline), imposta
\code{p\_semAdd}, incrementa \code{softBlockCount} \emph{solo} se il
semaforo è un semaforo di device reale (esclude quindi sia i semafori
``normali'' di sincronizzazione sia --- tramite \code{isDeviceSemaphore}
--- lo pseudo-clock, il cui incremento è gestito esplicitamente dal
chiamante quando serve, es. in \code{CLOCKWAIT}), inserisce il processo
nella ASL con \code{insertBlocked}, azzera \code{currentProcess} e richiama
lo scheduler. \`E la funzione invocata da \code{PASSEREN} e da \code{DOIO}
quando il processo deve effettivamente bloccarsi.

\begin{lstlisting}
static void activeProcs_add(pcb_t *p) {
    for (int i = 0; i < MAXPROC; i++) {
        if (activeProcs[i] == NULL) { activeProcs[i] = p; return; }
    }
    PANIC();
}

static void activeProcs_remove(pcb_t *p) {
    for (int i = 0; i < MAXPROC; i++) {
        if (activeProcs[i] == p) { activeProcs[i] = NULL; return; }
    }
}

static pcb_t *findProcessByPid(int target) {
    for (int i = 0; i < MAXPROC; i++) {
        if (activeProcs[i] != NULL && activeProcs[i]->p_pid == target)
            return activeProcs[i];
    }
    return NULL;
}
\end{lstlisting}
Tre funzioni di gestione dell'array \code{activeProcs}, tutte $O(\text{MAXPROC})$
per scansione lineare:
\begin{itemize}
\item \func{activeProcs\_add(p)}: trova il primo slot libero
  (\code{NULL}) e vi inserisce \code{p}. Se non trova spazio (impossibile in
  condizioni normali, dato che \code{allocPcb} fallirebbe prima se i 20 PCB
  sono esauriti) chiama \code{PANIC()} per difesa.
\item \func{activeProcs\_remove(p)}: trova lo slot che contiene esattamente
  \code{p} e lo azzera.
\item \func{findProcessByPid(target)}: scansione lineare per trovare il
  processo (vivo, non necessariamente pronto) con un dato PID; ritorna
  \code{NULL} se non esiste (PID già terminato o mai esistito). Usata da
  \code{TERMPROCESS} quando il target non è il processo corrente.
\end{itemize}

\subsection{\func{terminateProcess} --- terminazione ricorsiva}
\begin{lstlisting}
static void terminateProcess(pcb_t *p) {
    if (p == NULL) return;
    if (p->p_pid == -1) return;

    p->p_pid = -1;

    pcb_t *child;
    while ((child = removeChild(p)) != NULL) {
        terminateProcess(child);
    }

    outChild(p);
    activeProcs_remove(p);

    if (p->p_semAdd != NULL) {
        int *sem = p->p_semAdd;
        pcb_t *removed = outBlocked(p);
        p->p_semAdd = NULL;
        if (removed != NULL) {
            if (isDeviceSemaphore(sem) || sem == &devSems[PSEUDOCLK_SEM]) {
                softBlockCount--;
            } else {
                (*sem)++;
            }
        }
    }

    outProcQ(&readyQueue, p);

    if (p == currentProcess) {
        currentProcess = NULL;
    }

    freePcb(p);
    processCount--;
}
\end{lstlisting}
Implementa la semantica di \code{TERMPROCESS} (SYS2): terminare un processo
significa terminare \emph{ricorsivamente tutto il sotto-albero} radicato in
esso (i figli, i figli dei figli, ecc.), come richiesto dalla specifica di
Pandos. Analisi:

\begin{enumerate}
\item \textbf{Guardia contro doppie terminazioni}: se \code{p} è
  \code{NULL} o ha già \code{p\_pid == -1} (segnale interno ``sto già
  terminando/terminato''), ritorna subito. Questo flag protegge da eventuali
  riferimenti circolari o rientri accidentali durante la ricorsione.
\item \code{p->p\_pid = -1}: marca \code{p} come ``in fase di
  terminazione'' \emph{prima} di ricorrere sui figli, per evitare che una
  struttura anomala porti a rivisitarlo.
\item \textbf{Ricorsione sui figli}: il ciclo \code{while
  ((child=removeChild(p)) != NULL)} stacca ripetutamente il \emph{primo}
  figlio di \code{p} e lo termina ricorsivamente; poiché \code{removeChild}
  accorcia la lista dei figli a ogni chiamata, il ciclo termina quando
  \code{p} non ha più figli. Questa è una visita in stile
  \emph{depth-first}: prima si smaltisce interamente un ramo (e i suoi
  discendenti), poi si passa al fratello successivo.
\item \code{outChild(p)}: stacca \code{p} stesso dalla lista dei figli del
  \emph{proprio} padre. Necessario perché, come spiegato in Fase 1, senza
  questo passo un processo che termina se stesso (o che viene terminato
  senza che il padre stia già iterando su di lui) resterebbe agganciato come
  ``figlio fantasma''.
\item \code{activeProcs\_remove(p)}: toglie \code{p} dall'array globale dei
  processi vivi.
\item \textbf{Se \code{p} era bloccato su un semaforo}
  (\code{p->p\_semAdd != NULL}): va tolto dalla coda di quel semaforo
  (\code{outBlocked}) e --- punto delicato --- bisogna decidere \emph{cosa
  fare del valore del semaforo}:
  \begin{itemize}
  \item se è un semaforo di device \emph{oppure} lo pseudo-clock: il
    processo era ``soft-blocked'', cioè aspettava un evento che
    \emph{arriverà comunque} dall'hardware (l'interrupt del device o il
    prossimo tick). Il valore del semaforo \textbf{non} va toccato: sarà
    l'interrupt (quando arriverà) a fare la sua V regolare. Si aggiorna
    solo il contatore \code{softBlockCount--} (un processo soft-blocked in
    meno).
  \item altrimenti (semaforo ``normale'', di sincronizzazione tra
    processi): la \code{P} che il processo terminato aveva eseguito
    \emph{va annullata} incrementando il semaforo (\code{(*sem)++}),
    altrimenti il suo valore resterebbe sbilanciato verso il basso e un
    futuro processo che tenta la \code{P} si bloccherebbe per sempre
    aspettando un rilascio che non arriverà mai (il processo che avrebbe
    dovuto fare la V corrispondente non esiste più).
  \end{itemize}
\item \code{outProcQ(\&readyQueue, p)}: rimozione difensiva dalla ready
  queue (se \code{p} non vi si trova, la funzione non fa nulla e ritorna
  \code{NULL} senza errori: si veda Fase 1).
\item Se \code{p} era \code{currentProcess}, lo si azzera esplicitamente:
  il chiamante (\code{TERMPROCESS} in \code{syscallHandler}) userà questo
  segnale per capire se deve richiamare direttamente lo scheduler.
\item \code{freePcb(p)}: il PCB torna alla free list della Fase 1.
\item \code{processCount--}: un processo vivo in meno.
\end{enumerate}

\subsection{\func{uTLB\_RefillHandler} --- TLB-Refill event}
\begin{lstlisting}
void uTLB_RefillHandler(void) {
    state_t *savedState = (state_t *) BIOSDATAPAGE;

#ifdef SUPPORT_LEVEL
    unsigned int vpn = (savedState->entry_hi >> VPNSHIFT) & 0xFFFFF;
    int index;
    if (vpn == 0xBFFFF) index = 31;
    else index = (int)(vpn - 0x80000);

    pteEntry_t *pte = &currentProcess->p_supportStruct->sup_privatePgTbl[index];
    setENTRYHI(pte->pte_entryHI);
    setENTRYLO(pte->pte_entryLO);
    TLBWR();
    LDST(savedState);
#else
    setENTRYHI(0x80000000);
    setENTRYLO(0x00000000);
    TLBWR();
    LDST(savedState);
#endif
}
\end{lstlisting}
Compilato \emph{condizionalmente} in base alla macro \code{SUPPORT\_LEVEL}
(definita solo quando si compila il target \code{phase3}). Nella Fase 2
pura non esiste memoria virtuale a livello utente: il ramo \code{\#else} è
solo uno skeleton placeholder ereditato dal tester (scrive
un'entry fissa per la pagina di test \code{p5}), mai realmente esercitato
in modo significativo. Nella Fase 3 (\code{\#ifdef SUPPORT\_LEVEL}) è il
\textbf{percorso veloce} dei page fault: quando la traduzione manca dal TLB
\emph{ma è già presente} nella Page Table privata della U-proc
(\code{currentProcess->p\_supportStruct->sup\_privatePgTbl}), questo
handler si limita a copiarla dentro il TLB (\code{setENTRYHI}/\code{
setENTRYLO}/\code{TLBWR}) e a riprendere l'esecuzione (\code{LDST}), senza
mai coinvolgere il Pager né il disco. La discussione approfondita di questo
meccanismo (e della distinzione rispetto al Pager) è nel documento di Fase
3.

\subsection{\func{tlbExceptionHandler} e \func{programTrapHandler}}
\begin{lstlisting}
static void tlbExceptionHandler(void) {
    passUpOrDie(PGFAULTEXCEPT);
}

static void programTrapHandler(void) {
    passUpOrDie(GENERALEXCEPT);
}
\end{lstlisting}
Due semplici wrapper: qualunque eccezione di gestione memoria (page fault)
viene delegata a \code{passUpOrDie} con tipo \code{PGFAULTEXCEPT} (0);
qualunque altra eccezione non gestita direttamente dal Nucleus (istruzione
illegale, misallineamento, ecc.) con tipo \code{GENERALEXCEPT} (1). Questi
due indici corrispondono esattamente agli slot \code{[0]}/\code{[1]} degli
array \code{sup\_exceptState}/\code{sup\_exceptContext} nella
\code{support\_t} (Fase 1/3): è così che il Nucleus, senza sapere nulla del
Support Level, riesce comunque a smistare l'eccezione al gestore corretto di
livello superiore.

\subsection{\func{exceptionHandler} --- il dispatcher principale}
\begin{lstlisting}
void exceptionHandler(void) {
    state_t *savedState = (state_t *) BIOSDATAPAGE;
    unsigned int cause   = savedState->cause;

    updateCPUTime();

    if (cause & 0x80000000) {
        interruptHandler();
        return;
    }

    unsigned int excCode = cause & CAUSE_EXCCODE_MASK;

    switch (excCode) {
        case EXC_ECU:
        case EXC_ECM:
            savedState->pc_epc += WORDLEN;
            syscallHandler(savedState);
            break;

        case EXC_IPF:
        case EXC_LPF:
        case EXC_SPF:
        case EXC_MOD:
        case EXC_TLBL:
        case EXC_TLBS:
        case EXC_UTLBL:
        case EXC_UTLBS:
            tlbExceptionHandler();
            break;

        default:
            programTrapHandler();
            break;
    }
}
\end{lstlisting}
È l'\textbf{unico punto di ingresso} registrato nel Pass Up Vector per tutte
le eccezioni non-refill. Sequenza:
\begin{enumerate}
\item Legge \code{savedState} da \code{BIOSDATAPAGE} e il registro
  \code{cause} salvato dal firmware.
\item \code{updateCPUTime()}: aggiorna \emph{subito} la contabilità del
  tempo CPU del processo appena interrotto, prima di fare qualunque altra
  cosa --- così anche il tempo trascorso mentre si decide cosa fare
  dell'eccezione viene attribuito correttamente.
\item \textbf{Bit 31 di \code{cause} acceso $\to$ interrupt}: delega
  immediatamente a \code{interruptHandler()} (Fase 2, modulo
  \code{interrupts.c}) e ritorna (l'interrupt handler non ritornerà mai
  qui: fa sempre \code{LDST} o chiama lo scheduler).
\item Altrimenti, legge il campo basso \code{excCode} e smista:
  \begin{itemize}
  \item \code{EXC\_ECU} (8, \code{ecall} da user-mode) o \code{EXC\_ECM}
    (11, \code{ecall} da kernel-mode, usata dal Nucleus stesso e dai
    processi kernel come \code{test}): \textbf{avanza il PC di una word}
    (\code{WORDLEN=4}) \emph{prima} di gestire la syscall. Questo è
    essenziale: senza avanzare il PC, quando il processo verrà ripreso
    rieseguirebbe la stessa istruzione \code{ecall} all'infinito. Poi
    chiama \code{syscallHandler(savedState)}.
  \item un insieme di codici legati a page fault/TLB (\code{EXC\_IPF},
    \code{EXC\_LPF}, \code{EXC\_SPF} --- i codici RISC-V standard per
    Instruction/Load/Store Page Fault --- più i codici specifici
    dell'architettura $\mu$RISCV \code{EXC\_MOD}, \code{EXC\_TLBL},
    \code{EXC\_TLBS}, \code{EXC\_UTLBL}, \code{EXC\_UTLBS}): tutti trattati
    uniformemente come ``memory management trap'' $\to$
    \code{tlbExceptionHandler()}.
  \item \code{default}: qualunque altra causa (istruzione illegale,
    accesso non allineato, breakpoint, ecc.) $\to$
    \code{programTrapHandler()}.
  \end{itemize}
\end{enumerate}

\subsection{\func{syscallHandler} --- le 10 syscall del Nucleus}
\begin{lstlisting}
static void syscallHandler(state_t *savedState) {
    int sysCode = (int) savedState->reg_a0;

    if ((savedState->status & MSTATUS_MPP_MASK) == 0 && sysCode < 0) {
        savedState->cause = PRIVINSTR;
        programTrapHandler();
        return;
    }

    if (sysCode >= 1) {
        passUpOrDie(GENERALEXCEPT);
        return;
    }

    switch (sysCode) {
        case CREATEPROCESS: { ... }
        case TERMPROCESS:   { ... }
        case PASSEREN:      { ... }
        case VERHOGEN:      { ... }
        case DOIO:          { ... }
        case GETTIME:       { ... }
        case CLOCKWAIT:     { ... }
        case GETSUPPORTPTR: { ... }
        case GETPROCESSID:  { ... }
        case YIELD:         { ... }
        default:
            passUpOrDie(GENERALEXCEPT);
    }
}
\end{lstlisting}

\subsubsection{Protezione dei privilegi}
Il numero di syscall arriva in \code{a0} (convenzione RISC-V:
\code{reg\_a0}). Le syscall del Nucleus sono numerate con valori
\textbf{negativi} da \code{-1} (\code{CREATEPROCESS}) a \code{-10}
(\code{YIELD}); i valori \textbf{positivi} sono riservati alle syscall del
Support Level (Fase 3). Prima regola:
\begin{quote}
\code{if ((savedState->status \& MSTATUS\_MPP\_MASK) == 0 \&\& sysCode < 0)}
\end{quote}
Se il campo \code{MPP} dello \code{status} salvato è 0 (cioè il chiamante
\emph{era} in \textbf{user-mode}) \emph{e} il codice richiesto è negativo
(una syscall riservata al kernel), l'eccezione viene forzatamente
riqualificata come \code{PRIVINSTR} (istruzione privilegiata, tentativo di
uso non autorizzato) e delegata a \code{programTrapHandler()}: nessuna
U-proc può invocare direttamente una syscall del Nucleus.

Seconda regola: se \code{sysCode >= 1} (syscall del Support Level, di
competenza di un livello \emph{superiore} al Nucleus, che non sa nulla di
come gestirle), viene delegata con \code{passUpOrDie(GENERALEXCEPT)}: il
Nucleus non implementa mai syscall positive, le passa sempre su.

\subsubsection{SYS1 --- \code{CREATEPROCESS}}
\begin{lstlisting}
case CREATEPROCESS: {
    state_t   *newState = (state_t *) savedState->reg_a1;
    int        prio     = (int) savedState->reg_a2;
    support_t *support  = (support_t *) savedState->reg_a3;

    pcb_t *child = allocPcb();
    if (!child) {
        savedState->reg_a0 = (unsigned int) -1;
        copyState(&currentProcess->p_s, savedState);
        LDST(&currentProcess->p_s);
    }

    copyState(&child->p_s, newState);
    child->p_supportStruct = support;
    child->p_prio          = prio;

    activeProcs_add(child);
    insertChild(currentProcess, child);
    processCount++;

    savedState->reg_a0 = (unsigned int) child->p_pid;
    copyState(&currentProcess->p_s, savedState);
    insertProcQ(&readyQueue, currentProcess);
    insertProcQ(&readyQueue, child);
    currentProcess = NULL;
    scheduler();
    break;
}
\end{lstlisting}
Crea un nuovo processo figlio del chiamante. Argomenti (convenzione: \code{a1},
\code{a2}, \code{a3}): puntatore allo stato iniziale del processore che il
nuovo processo dovrà avere, priorità, puntatore opzionale a una
\code{support\_t}. Passi:
\begin{enumerate}
\item \code{allocPcb()}: se la free list dei PCB è esaurita (limite di 20
  processi vivi raggiunto), \code{a0} del \emph{chiamante} viene impostato a
  \code{-1} (fallimento) e si riprende immediatamente il chiamante con
  \code{LDST} --- \emph{niente} nuovo processo, nessuna modifica a
  \code{processCount}.
\item Altrimenti si copia lo stato iniziale fornito nel PCB del figlio
  (\code{copyState}), si registrano support struct e priorità.
\item Il figlio viene aggiunto ad \code{activeProcs} e inserito nell'albero
  come figlio di \code{currentProcess} (\code{insertChild}).
  \code{processCount++}.
\item \code{a0} del chiamante viene impostato al \textbf{PID del figlio}
  appena creato (valore di ritorno della syscall lato padre).
\item \textbf{Sia il padre sia il figlio} vengono inseriti nella ready
  queue (in ordine: prima il padre, poi il figlio, ma \code{insertProcQ}
  riordina comunque per priorità), e si richiama lo scheduler. Questo
  \emph{abilita la preemption immediata}: se il figlio ha priorità
  superiore al padre, sarà lui il primo a essere estratto ed eseguito,
  anche prima che il padre riprenda dal punto della \code{ecall}.
\end{enumerate}

\subsubsection{SYS2 --- \code{TERMPROCESS}}
\label{subsub:sys2-termprocess}
\begin{lstlisting}
case TERMPROCESS: {
    int targetPid = (int) savedState->reg_a1;
    updateCPUTime();

    pcb_t *target = (targetPid == 0) ? currentProcess : findProcessByPid(targetPid);
    if (!target) {
        LDST(savedState);
    } else {
        int terminatingSelf = (target == currentProcess);
        terminateProcess(target);
        if (terminatingSelf) currentProcess = NULL;
    }

    if (currentProcess == NULL) {
        scheduler();
    } else {
        copyState(&currentProcess->p_s, savedState);
        LDST(&currentProcess->p_s);
    }
    break;
}
\end{lstlisting}
Se \code{a1 == 0} il target è il processo chiamante stesso; altrimenti si
cerca per PID con \code{findProcessByPid}. Se il target non esiste più
(magari già terminato da un altro processo), si riprende semplicemente il
chiamante senza fare nulla (\code{LDST(savedState)}: nota che qui si usa
direttamente \code{savedState}, non \code{\&currentProcess->p\_s}, perché il
chiamante --- che \emph{non} è il target in questo ramo --- non ha ancora
subito modifiche al proprio stato salvato). Se il target esiste, si chiama
\code{terminateProcess} (che termina ricorsivamente lui e tutto il suo
sotto-albero, si veda sopra). Se il chiamante ha terminato \emph{se stesso},
\code{currentProcess} è già stato azzerato da \code{terminateProcess}, e si
richiama direttamente lo scheduler; altrimenti (il chiamante ha terminato
\emph{un altro} processo, es. un fratello o un discendente) il chiamante
stesso prosegue: si salva il suo stato e si riprende con \code{LDST}.

\subsubsection{SYS3 --- \code{PASSEREN} (P)}
\begin{lstlisting}
case PASSEREN: {
    int *semAddr = (int *) savedState->reg_a1;
    updateCPUTime();
    copyState(&currentProcess->p_s, savedState);
    (*semAddr)--;
    if (*semAddr < 0) {
        blockCurrentProcess(semAddr);
    }
    LDST(&currentProcess->p_s);
    break;
}
\end{lstlisting}
Implementa la semantica classica del semaforo a \textbf{contatore
negativo}: si decrementa sempre il valore; se il risultato è negativo
significa che non c'era disponibilità, e il processo si blocca
(\code{blockCurrentProcess}, che a sua volta chiama \code{scheduler()} e non
ritorna qui). Se invece il valore resta $\ge 0$, il processo continua
subito: si riprende con \code{LDST}. Da notare che lo stato viene salvato
nel PCB (\code{copyState}) \emph{prima} di decidere se bloccarsi, cosicché
se il blocco avviene, \code{blockCurrentProcess} trova già uno stato
coerente salvato in \code{p\_s}.

\subsubsection{SYS4 --- \code{VERHOGEN} (V)}
\begin{lstlisting}
case VERHOGEN: {
    int *semAddr = (int *) savedState->reg_a1;
    updateCPUTime();
    copyState(&currentProcess->p_s, savedState);
    (*semAddr)++;
    if (*semAddr <= 0) {
        pcb_t *unblocked = removeBlocked(semAddr);
        if (unblocked) {
            unblocked->p_semAdd = NULL;
            insertProcQ(&readyQueue, unblocked);
        }
    }
    LDST(&currentProcess->p_s);
    break;
}
\end{lstlisting}
Incrementa il semaforo; se il valore risultante resta $\le 0$, significa che
c'era (almeno) un processo in coda: lo si rimuove dalla ASL
(\code{removeBlocked}, che restituisce il primo in coda in ordine FIFO) e lo
si reinserisce nella ready queue. Il chiamante della V \emph{non si blocca
mai}: dopo l'eventuale sblocco altrui, riprende sempre subito con
\code{LDST}.

\subsubsection{SYS5 --- \code{DOIO}}
\begin{lstlisting}
case DOIO: {
    int *commandAddr  = (int *) savedState->reg_a1;
    int  commandValue = (int) savedState->reg_a2;

    unsigned int offset = (unsigned int)commandAddr - START_DEVREG;
    int line    = (int)(offset / 0x80) + 3;
    int dev     = (int)((offset % 0x80) / 0x10);
    int subword = (int)((offset % 0x10) / WORDLEN);

    int semIdx = (line == 7)
        ? ((subword == 3) ? TERM_TX_SEM(dev) : TERM_RX_SEM(dev))
        : DEV_SEM_BASE(line, dev);

    updateCPUTime();
    copyState(&currentProcess->p_s, savedState);

    *commandAddr = commandValue;
    devSems[semIdx]--;
    blockCurrentProcess(&devSems[semIdx]);
    break;
}
\end{lstlisting}
Avvia un'operazione di I/O e blocca \emph{sempre} il processo chiamante,
perché ogni operazione di I/O è per natura asincrona e si conclude con un
interrupt di completamento gestito da \code{interrupts.c}.
\begin{enumerate}
\item Calcolo del \textbf{semaforo giusto} a partire dall'indirizzo del
  registro di comando passato dal chiamante (\code{commandAddr}):
  sottraendo la base \code{START\_DEVREG} si ottiene un offset, da cui si
  ricavano linea di interrupt (\code{line}, ogni linea occupa \code{0x80}
  byte), numero di device sulla linea (\code{dev}, ogni device occupa
  \code{0x10} byte) e la ``subword'' dentro il blocco di registri del
  device (\code{subword}, ogni registro è una word da 4 byte). Per la linea
  terminale (\code{line == 7}), che ha 4 registri per device (status/comando
  di ricezione, poi status/comando di trasmissione), la subword \code{3}
  identifica il registro di comando \emph{trasmissione} $\to$ semaforo TX,
  altrimenti si usa il semaforo RX.
\item Si aggiorna la contabilità del tempo CPU e si salva lo stato corrente.
\item \textbf{Si scrive il comando sul registro del device}
  (\code{*commandAddr = commandValue}): questo è l'istante in cui l'I/O
  viene effettivamente avviato a livello hardware.
\item Si decrementa il semaforo del device e si blocca il processo su di
  esso (\code{blockCurrentProcess}).
\end{enumerate}
\textbf{Assenza di race condition}: tra la scrittura del comando e il blocco
del processo il Nucleus gira sempre con gli interrupt disabilitati (il
Nucleus in generale processa le eccezioni con interrupt mascherati fino al
prossimo \code{WAIT}/\code{LDST}); anche se il device completasse
l'operazione ``istantaneamente'', l'interrupt corrispondente resterebbe
\emph{pending} e verrebbe servito solo dopo che lo scheduler ha eseguito
\code{WAIT} o è ripartito con un altro processo --- non si perde mai il
segnale di completamento.

\subsubsection{SYS6 --- \code{GETTIME}}
\begin{lstlisting}
case GETTIME: {
    updateCPUTime();
    savedState->reg_a0 = (unsigned int) currentProcess->p_time;
    LDST(savedState);
    break;
}
\end{lstlisting}
Aggiorna la contabilità (per includere anche il tempo trascorso fino a
questo istante nella syscall stessa) e restituisce in \code{a0} il tempo
CPU totale accumulato dal processo chiamante. Nessun blocco: si riprende
subito con \code{LDST(savedState)} (usa direttamente lo stato originale
letto da \code{BIOSDATAPAGE}, dato che non serve passare per il PCB in
questo caso semplice).

\subsubsection{SYS7 --- \code{CLOCKWAIT}}
\begin{lstlisting}
case CLOCKWAIT: {
    updateCPUTime();
    copyState(&currentProcess->p_s, savedState);
    devSems[PSEUDOCLK_SEM]--;
    softBlockCount++;
    blockCurrentProcess(&devSems[PSEUDOCLK_SEM]);
    break;
}
\end{lstlisting}
Blocca il processo chiamante fino al prossimo tick dello pseudo-clock (ogni
\code{PSECOND}, 100\,ms). A differenza di \code{PASSEREN} generica, qui
\code{softBlockCount++} viene incrementato \textbf{esplicitamente} \emph{prima}
di chiamare \code{blockCurrentProcess}: dato che il semaforo pseudo-clock è
per definizione escluso da \code{isDeviceSemaphore} (per evitare doppio
conteggio quando \emph{altri} punti del codice bloccano su semafori
generici), qui il chiamante si assume esplicitamente la responsabilità di
incrementare il contatore, così che lo scheduler possa comunque distinguere
correttamente ``sto aspettando il prossimo tick'' da un deadlock.

\subsubsection{SYS8 --- \code{GETSUPPORTPTR}}
\begin{lstlisting}
case GETSUPPORTPTR: {
    savedState->reg_a0 = (unsigned int) currentProcess->p_supportStruct;
    LDST(savedState);
    break;
}
\end{lstlisting}
Restituisce semplicemente il puntatore alla \code{support\_t} del processo
chiamante (\code{NULL} se non ne ha una). È la syscall con cui, in Fase 3,
il Pager e il general exception handler recuperano la propria support
structure appena vengono attivati dal pass-up.

\subsubsection{SYS9 --- \code{GETPROCESSID}}
\begin{lstlisting}
case GETPROCESSID: {
    int wantParent = (int) savedState->reg_a1;
    savedState->reg_a0 = (wantParent == 0)
        ? (unsigned int) currentProcess->p_pid
        : (currentProcess->p_parent ? (unsigned int) currentProcess->p_parent->p_pid : 0);
    LDST(savedState);
    break;
}
\end{lstlisting}
Se \code{a1 == 0} restituisce il PID del chiamante stesso; se \code{a1 !=
0} restituisce il PID del \emph{padre} del chiamante, oppure \code{0} se il
chiamante non ha un padre (è il processo radice).

\subsubsection{SYS10 --- \code{YIELD}}
\begin{lstlisting}
case YIELD: {
    updateCPUTime();
    copyState(&currentProcess->p_s, savedState);
    yieldedProcess = currentProcess;
    currentProcess = NULL;
    scheduler();
    break;
}
\end{lstlisting}
Il processo chiamante cede volontariamente la CPU. \emph{Non} viene
reinserito subito nella ready queue: viene invece registrato in
\code{yieldedProcess}, e sarà lo scheduler (passo 0, discusso sopra) a
decidere se rimetterlo in coda (perché esiste un processo pronto di
priorità pari o superiore) oppure farlo ripartire immediatamente con una
nuova time slice (se è comunque il candidato più prioritario). Questo
meccanismo indiretto garantisce che uno \code{YIELD} \emph{ceda davvero} la
CPU quando ha senso farlo, invece di rientrare subito in
\code{insertProcQ} e ripescarsi da solo in coda per un caso limite.

\subsection{\func{passUpOrDie} --- il meccanismo di pass-up-or-die}
\begin{lstlisting}
static void passUpOrDie(int exceptionType) {
    if (!currentProcess || !currentProcess->p_supportStruct) {
        updateCPUTime();
        terminateProcess(currentProcess);
        currentProcess = NULL;
        scheduler();
    } else {
        support_t *sup = currentProcess->p_supportStruct;
        copyState(&sup->sup_exceptState[exceptionType], (state_t *) BIOSDATAPAGE);
        context_t *ctx = &sup->sup_exceptContext[exceptionType];
        LDCXT(ctx->stackPtr, ctx->status, ctx->pc);
    }
}
\end{lstlisting}
Il cuore del disaccoppiamento tra Nucleus e Support Level. Chiamato ogni
volta che il Nucleus incontra un'eccezione che \emph{non sa} gestire
direttamente (page fault, program trap, o syscall con codice $\ge 1$):
\begin{itemize}
\item \textbf{Caso ``die''}: se non c'è un processo corrente, oppure il
  processo corrente \emph{non ha} una support structure
  (\code{p\_supportStruct == NULL}, tipico di un processo puramente Fase 2
  come \code{test} in modalità pura, che non ha nessun livello superiore a
  cui rivolgersi), il processo viene semplicemente terminato con
  \code{terminateProcess} (che, ricordiamo, termina anche tutto il suo
  sotto-albero), e si richiama lo scheduler.
\item \textbf{Caso ``pass up''}: se invece il processo \emph{ha} una
  support structure (una U-proc di Fase 3), lo stato salvato in
  \code{BIOSDATAPAGE} viene copiato nello slot corrispondente
  (\code{sup\_exceptState[PGFAULTEXCEPT]} o
  \code{sup\_exceptState[GENERALEXCEPT]}, a seconda del tipo passato come
  argomento), e si esegue un \code{LDCXT} con lo stack pointer, lo status e
  il PC salvati nel \code{context\_t} corrispondente
  (\code{sup\_exceptContext[...]}): è un vero salto (non-locale) a un
  gestore di livello Support (il Pager o il general exception handler,
  configurati in Fase 3 da \code{launchUproc}), su uno stack dedicato
  interamente separato da quello del Nucleus.
\end{itemize}
Questo è esattamente ciò che permette alla Fase 3 di essere costruita
\emph{sopra} il Nucleus senza toccarne il codice: il Nucleus si limita a
``spedire'' l'eccezione al livello giusto, senza sapere nulla di cosa
significhino \code{PGFAULTEXCEPT}/\code{GENERALEXCEPT} per chi le riceve.

% ============================================================
\section{\code{phase2/interrupts.c} --- gestione degli interrupt}
% ============================================================

\subsection{Costanti e utility}
\begin{lstlisting}
#define INT_BITMAP_BASE  0x10000040
#define INT_BITMAP(line) (*((unsigned int *)(INT_BITMAP_BASE + ((line) - 3) * 0x4)))

#define DEV_REG_BASE(line, dev) \
    ((unsigned int *)(START_DEVREG + ((line) - 3) * 0x80 + (dev) * 0x10))

#define DEV_STATUS(base)          ((base)[0])
#define DEV_COMMAND(base)         ((base)[1])

#define TERM_RECV_STATUS(base)    ((base)[0])
#define TERM_RECV_COMMAND(base)   ((base)[1])
#define TERM_TRANSM_STATUS(base)  ((base)[2])
#define TERM_TRANSM_COMMAND(base) ((base)[3])

static int getHighestPriorityDevice(unsigned int bitmap) {
    for (int i = 0; i < 8; i++) {
        if (bitmap & (1u << i)) return i;
    }
    return -1;
}
\end{lstlisting}
\code{INT\_BITMAP(line)} legge, per una data linea di interrupt (3--7), la
bitmap hardware a 8 bit che indica quali degli 8 device su quella linea
hanno un interrupt pendente. \code{DEV\_REG\_BASE(line, dev)} calcola
l'indirizzo del blocco di registri di un device specifico (ogni linea
occupa \code{0x80} byte, ogni device \code{0x10} byte all'interno della
linea). Le macro \code{DEV\_STATUS}/\code{DEV\_COMMAND} e le quattro
\code{TERM\_*} sono semplici alias per interpretare quel blocco come
registri (status, comando) di un device generico o, per i terminali, come
due coppie indipendenti (ricezione, trasmissione).
\func{getHighestPriorityDevice(bitmap)} scansiona i bit da 0 a 7 e ritorna
l'indice del \emph{primo} bit acceso (cioè il device a \textbf{priorità più
alta}, dato che per convenzione hardware il bit meno significativo ha
priorità maggiore); \code{-1} se nessun bit è acceso (interrupt spurio).

\subsection{\func{interruptHandler} --- struttura generale}
\begin{lstlisting}
void interruptHandler(void) {
    state_t *savedState = (state_t *) BIOSDATAPAGE;
    unsigned int cause   = savedState->cause;
    unsigned int excCode = cause & CAUSE_EXCCODE_MASK;

    cpu_t now;
    STCK(now);
    if (currentProcess != NULL) {
        currentProcess->p_time += (now - startTOD);
    }
    startTOD = now;

    if (excCode == 7u) { /* PLT */ ... }
    if (excCode == 3u) { /* Interval Timer */ ... }
    if (excCode >= 17u && excCode <= 21u) { /* device */ ... }

    if (currentProcess != NULL) LDST(savedState);
    else scheduler();
}
\end{lstlisting}
Legge \code{cause} e il suo \code{excCode}, aggiorna la contabilità del
tempo CPU (stessa logica di \code{updateCPUTime}, qui reimplementata
inline), e poi distingue tre grandi classi in base al valore di
\code{excCode}: \code{7} (PLT), \code{3} (Interval Timer),
\code{17--21} (le cinque linee hardware dei device esterni, che nel
registro \code{cause} sono mappate a partire da \code{17} pur
corrispondendo alle linee fisiche \code{3--7}, da cui il ``\code{-14}''
visto più avanti). Se nessuno di questi rami restituisce prima (return),
il codice finale ripristina \code{currentProcess} con \code{LDST} se
esiste ancora, altrimenti chiama lo scheduler.

\subsection{PLT --- fine time slice (preemption)}
\begin{lstlisting}
if (excCode == 7u) {
    setTIMER((cpu_t) NEVER);
    if (currentProcess != NULL) {
        copyState(&currentProcess->p_s, savedState);
        currentProcess->p_s.status |= MSTATUS_MIE_MASK;
        insertProcQ(&readyQueue, currentProcess);
        currentProcess = NULL;
    }
    scheduler();
    return;
}
\end{lstlisting}
Questo è il meccanismo che realizza la \textbf{preemption}: quando la time
slice di 5\,ms del processo corrente scade, il PLT genera questo interrupt.
Il gestore:
\begin{enumerate}
\item \textbf{Disarma il PLT} caricandogli il valore massimo possibile
  (\code{NEVER = 0x7FFFFFFF}), così non genererà ulteriori interrupt finché
  lo scheduler non ne programma uno nuovo per il prossimo processo.
\item Se c'era un processo in esecuzione, ne salva lo stato
  (\code{copyState}), \textbf{si assicura che, alla prossima ripartenza,
  abbia gli interrupt abilitati} (forzando il bit \code{MSTATUS\_MIE\_MASK}
  nello \code{status} salvato --- precauzione difensiva, dato che lo stato
  salvato dal firmware al momento dell'interrupt riflette lo status
  \emph{precedente} all'interrupt stesso), e lo rimette nella ready queue
  (\code{insertProcQ}, che lo riposiziona secondo la sua priorità: un
  round-robin ``puro'' tra processi di pari priorità).
\item Richiama lo scheduler, che sceglierà il prossimo processo da
  eseguire (possibilmente lo stesso, se non c'è nessun altro pronto).
\end{enumerate}

\subsection{Interval Timer --- tick dello pseudo-clock}
\begin{lstlisting}
if (excCode == 3u) {
    LDIT(PSECOND);

    pcb_t *p;
    while ((p = removeBlocked(&devSems[PSEUDOCLK_SEM])) != NULL) {
        p->p_semAdd   = NULL;
        p->p_s.reg_a0 = 0;
        insertProcQ(&readyQueue, p);
        softBlockCount--;
    }

    devSems[PSEUDOCLK_SEM] = 0;

    if (currentProcess != NULL) LDST(savedState);
    else scheduler();
    return;
}
\end{lstlisting}
Ogni 100\,ms (\code{PSECOND}) l'Interval Timer genera questo interrupt.
\begin{enumerate}
\item \textbf{Riarma subito il timer} per il prossimo tick
  (\code{LDIT(PSECOND)}).
\item \textbf{Sblocca in blocco} tutti i processi che erano in attesa sul
  semaforo pseudo-clock (chiunque avesse chiamato \code{CLOCKWAIT}): il
  ciclo \code{while} continua a chiamare \code{removeBlocked} finché la
  coda non è vuota, e per ciascun processo estratto: azzera
  \code{p\_semAdd}, imposta \code{a0 = 0} (valore di ritorno di
  \code{CLOCKWAIT}), lo inserisce nella ready queue, e decrementa
  \code{softBlockCount}.
\item \textbf{Azzera esplicitamente} \code{devSems[PSEUDOCLK\_SEM]}: a
  differenza di un semaforo normale (dove ogni V incrementa di 1), qui si
  forza il valore a 0 perché lo pseudo-clock ha semantica ``broadcast'':
  \emph{tutti} i processi in attesa vengono liberati contemporaneamente a
  ogni tick, non uno alla volta come farebbe una sequenza di V singole.
\item Come sempre, se c'era un processo in esecuzione lo si riprende
  direttamente (\code{LDST(savedState)}, nessun cambio di processo
  necessario); altrimenti si chiama lo scheduler (che magari ora ha nuovi
  processi pronti, appena svegliati).
\end{enumerate}

\subsection{Interrupt di device (linee 3--7, \code{excCode} 17--21)}
\begin{lstlisting}
if (excCode >= 17u && excCode <= 21u) {
    int intLineNo = (int)(excCode - 14u);
    unsigned int bitmap = INT_BITMAP(intLineNo);
    int devNo = getHighestPriorityDevice(bitmap);

    if (devNo < 0) {
        if (currentProcess != NULL) LDST(savedState);
        else scheduler();
        return;
    }

    if (intLineNo == 7) {
        /* terminale: TX poi RX, vedi sotto */
    } else {
        unsigned int *devBase = DEV_REG_BASE(intLineNo, devNo);
        unsigned int savedStatus = DEV_STATUS(devBase);
        DEV_COMMAND(devBase) = ACK;

        int semIdx = DEV_SEM_BASE(intLineNo, devNo);
        if (devSems[semIdx] < 0) {
            devSems[semIdx]++;
            pcb_t *unblocked = removeBlocked(&devSems[semIdx]);
            if (unblocked != NULL) {
                unblocked->p_s.reg_a0 = savedStatus;
                unblocked->p_semAdd   = NULL;
                insertProcQ(&readyQueue, unblocked);
                softBlockCount--;
            }
        }
    }

    if (currentProcess != NULL) LDST(savedState);
    else scheduler();
    return;
}
\end{lstlisting}
\begin{enumerate}
\item \code{intLineNo = excCode - 14}: converte il codice \code{cause}
  (17--21) nella linea hardware effettiva (3--7: \code{IL\_DISK},
  \code{IL\_FLASH}, \code{IL\_ETHERNET}, \code{IL\_PRINTER},
  \code{IL\_TERMINAL}).
\item Legge la bitmap dei device con interrupt pendente su quella linea, e
  sceglie \code{devNo} col bit più basso acceso (priorità più alta). Se
  nessun bit è acceso, l'interrupt è spurio: si riprende/schedula senza
  fare nulla.
\item \textbf{Per device diversi dal terminale} (linee 3--6: disk, flash,
  ethernet, printer): si legge lo \textbf{status prima} di inviare l'ACK.
  Questo ordine è \emph{cruciale}: il registro di status del device viene
  tipicamente sovrascritto o azzerato nel momento in cui si invia il
  comando di ACK, quindi se si invertisse l'ordine il valore di status
  andrebbe perso e il processo risvegliato non saprebbe come è terminata la
  sua operazione di I/O. Dopo aver salvato lo status, si scrive
  \code{ACK} nel registro di comando (riconoscimento dell'interrupt), poi:
  se il semaforo del device era negativo (c'era almeno un processo in
  attesa), lo si incrementa e si sblocca il primo processo in coda
  (\code{removeBlocked}), impostandone \code{a0} allo status letto (valore
  di ritorno di \code{DOIO}) e decrementando \code{softBlockCount}. Se il
  semaforo \emph{non} era negativo (nessuno in attesa --- caso anomalo dato
  che \code{DOIO} blocca sempre, ma gestito difensivamente), non si fa
  nulla.
\item Al termine, come sempre, si riprende il processo corrente o si
  schedula.
\end{enumerate}

\subsubsection{Il caso speciale del terminale (linea 7)}
\begin{lstlisting}
unsigned int *termBase = DEV_REG_BASE(intLineNo, devNo);
unsigned int txStatus = TERM_TRANSM_STATUS(termBase) & 0xFFu;
unsigned int rxStatus = TERM_RECV_STATUS(termBase) & 0xFFu;

if (txStatus == OKCHARTRANS) {
    unsigned int savedStatus = TERM_TRANSM_STATUS(termBase);
    TERM_TRANSM_COMMAND(termBase) = ACK;
    int semIdx = TERM_TX_SEM(devNo);
    if (devSems[semIdx] < 0) {
        devSems[semIdx]++;
        pcb_t *unblocked = removeBlocked(&devSems[semIdx]);
        if (unblocked != NULL) {
            unblocked->p_s.reg_a0 = savedStatus;
            unblocked->p_semAdd   = NULL;
            list_add_tail(&unblocked->p_list, &readyQueue);
            softBlockCount--;
        }
    }
}

if (rxStatus == CHARRECV) {
    unsigned int savedStatus = TERM_RECV_STATUS(termBase);
    TERM_RECV_COMMAND(termBase) = ACK;
    int semIdx = TERM_RX_SEM(devNo);
    if (devSems[semIdx] < 0) {
        devSems[semIdx]++;
        pcb_t *unblocked = removeBlocked(&devSems[semIdx]);
        if (unblocked != NULL) {
            unblocked->p_s.reg_a0 = savedStatus;
            unblocked->p_semAdd   = NULL;
            insertProcQ(&readyQueue, unblocked);
            softBlockCount--;
        }
    }
}
\end{lstlisting}
Un terminale è in realtà \textbf{due sotto-device indipendenti}:
trasmissione (TX, scrittura di caratteri) e ricezione (RX, lettura di
caratteri), ciascuno con il proprio status/comando e il proprio semaforo
dedicato (\code{TERM\_TX\_SEM}/\code{TERM\_RX\_SEM}). Un singolo interrupt
sulla linea 7 può quindi segnalare il completamento di \emph{una} delle due
operazioni (o, in teoria, di entrambe se coincidono): il codice controlla
\textbf{prima} lo status di trasmissione (\code{OKCHARTRANS = 5}: carattere
trasmesso con successo) e \textbf{poi} quello di ricezione
(\code{CHARRECV = 5}: carattere ricevuto), gestendo ciascuno in modo
indipendente con la stessa logica (salva status, ACK, eventuale sblocco del
processo in attesa) vista sopra per i device generici. \textbf{TX ha
priorità su RX}: se entrambi gli status indicano un evento completato nello
stesso interrupt, si processa prima la trasmissione. Questo ordine evita
che un carattere ricevuto ``scavalchi'' silenziosamente un carattere appena
trasmesso in situazioni di traffico bidirezionale rapido, riducendo il
rischio di perdita di caratteri.

(Nota implementativa minore: nel ramo TX lo sblocco usa
\code{list\_add\_tail(\&unblocked->p\_list, \&readyQueue)} anziché
\code{insertProcQ}: essendo la ready queue ordinata per priorità tramite
\code{p\_prio}, e i processi risvegliati da un device tipicamente non hanno
priorità speciali diverse da quella con cui erano stati creati, l'effetto
pratico coincide quasi sempre con un \code{insertProcQ}; il ramo RX usa
invece coerentemente \code{insertProcQ}.)

% ============================================================
\section{Il ciclo completo di una operazione di I/O --- riepilogo narrativo}
% ============================================================

\`E lo scenario più chiesto all'orale: sapere raccontare, passo per passo,
cosa succede quando un processo chiede un'operazione di I/O.

\begin{enumerate}
\item Un processo esegue \code{SYSCALL(DOIO, indirizzoComando, comando, 0)}.
  Questo genera un'eccezione \code{ecall}, che il firmware intercetta
  salvando lo stato in \code{BIOSDATAPAGE} e saltando a
  \code{exceptionHandler} (l'indirizzo registrato nel Pass Up Vector).
\item \code{exceptionHandler} riconosce l'\code{ExcCode} come \code{ECALL},
  avanza il PC di una word, e chiama \code{syscallHandler}.
\item \code{syscallHandler} riconosce \code{sysCode == DOIO}, calcola il
  semaforo giusto dall'indirizzo del registro di comando, \textbf{scrive il
  comando sul device} e poi chiama \code{blockCurrentProcess}: il processo
  viene inserito nella ASL sul semaforo del device e \code{currentProcess}
  torna \code{NULL}.
\item Si richiama \code{scheduler()}. Se non ci sono altri processi pronti
  ma qualcuno è soft-blocked (come il processo appena bloccato), lo
  scheduler esegue \code{WAIT()}: la CPU virtuale si ferma in attesa di un
  interrupt, con tutti gli interrupt (tranne il PLT) abilitati.
\item Il device, terminata l'operazione richiesta, genera un
  \textbf{interrupt} sulla propria linea. Il firmware salva lo stato e
  salta di nuovo a \code{exceptionHandler}, che lo riconosce (bit 31 di
  \code{cause}) e delega a \code{interruptHandler}.
\item \code{interruptHandler} legge la bitmap della linea, individua il
  device, \textbf{legge lo status prima di fare l'ACK}, incrementa il
  semaforo del device e sblocca il processo in attesa con
  \code{removeBlocked}, impostandogli \code{a0} allo status letto.
\item Il processo risvegliato torna nella ready queue. Alla sua prossima
  esecuzione (dispatch dello scheduler), riprenderà dall'istruzione
  \emph{successiva} alla \code{ecall} della \code{DOIO}, trovando in
  \code{a0} lo status dell'operazione appena conclusa.
\end{enumerate}
Nessuna race condition è possibile perché tra l'emissione del comando (passo
3) e il blocco effettivo del processo, il Nucleus sta ancora gestendo
un'eccezione con gli interrupt mascherati: anche se il device completasse
l'operazione istantaneamente, l'interrupt resterebbe pending e verrebbe
processato solo dopo che il controllo è tornato a un punto sicuro
(\code{WAIT} o la ripresa di un altro processo).

% ============================================================
\section{Il test driver ufficiale: \code{phase2/p2test.c}}
% ============================================================

Come \code{p1test.c} per la Fase 1, \code{p2test.c} \textbf{non è codice
originale del progetto}: è il test driver standard del Nucleus (compilato
solo nel target \code{phase2}; fornisce la funzione \code{test()} che
\code{initial.c::main} lancia come primo processo). È volutamente
``contorto'' (lo dice il commento stesso in testa al file, \emph{``This is
pretty convoluted code, so good luck!''}) perché il suo scopo è mettere
sotto stress \emph{ogni} syscall e meccanismo del Nucleus in un'unica
esecuzione. Non viene analizzato riga per riga (è quasi 800 righe,
in gran parte ripetitive), ma è utile sapere cosa dimostra ciascun
sotto-processo, perché il professore può chiedere ``spiega cosa succede
quando gira \code{p2test}''.

\begin{center}
\begin{longtable}{@{}llp{8.6cm}@{}}
\toprule
\textbf{Processo} & \textbf{Ruolo} & \textbf{Cosa dimostra} \\
\midrule
\code{test} (``p1'') & radice & Fa 3 V e 3 P consecutive su un semaforo di
  test verificandone il valore atteso, poi costruisce a mano gli stati
  iniziali di tutti gli altri processi (stack decrescenti, ciascuno nella
  propria pagina) e li crea in sequenza con \code{CREATEPROCESS}. \\
\code{p2} & sincronizzazione + tempo & Si blocca finché \code{test} non fa
  \code{VERHOGEN(sem\_startp2)}; verifica il proprio PID
  (\code{GETPROCESSID}); esegue 20 cicli di V/P sullo stesso semaforo
  controllando che torni sempre a 0; confronta \code{GETTIME} col TOD reale
  (\code{STCK}) per verificare che la contabilità del tempo CPU sia
  corretta e coerente (mai più del tempo reale trascorso). \\
\code{p3} & pseudo-clock & Cicla su \code{CLOCKWAIT} misurando che ogni
  risveglio avvenga dopo circa un intervallo di clock, e che anche
  durante l'attesa passiva il tempo CPU registrato (\code{GETTIME}) resti
  coerente. \\
\code{p4} & terminazione a cascata & Crea una \emph{seconda incarnazione di
  se stesso} (stesso codice, PCB diverso) e attende che sia sincronizzata;
  quando il padre \code{p4} termina (\code{TERMPROCESS} su se stesso), il
  Nucleus deve terminare a cascata anche il figlio: verifica
  \code{terminateProcess} sui sotto-alberi. \\
\code{p5} (con \code{support\_t}) & pass-up-or-die \emph{con} support
  struct & È l'unico processo creato con una \code{support\_t} valida
  (\code{pFiveSupport}) anche in Fase 2 pura (nessuna vera memoria
  virtuale: serve solo a esercitare il meccanismo generico). Genera
  deliberatamente un accesso a un indirizzo invalido (program trap $\to$
  \code{p5gen}), poi un accesso fuori \code{kuseg} da user-mode (address
  error $\to$ ancora \code{p5gen}, che passa a gestione kernel-mode), poi
  una syscall dal Nucleus (\code{p5sys}) e un vero page fault
  (\code{p5mm}). Dimostra che \code{passUpOrDie} instrada correttamente le
  diverse eccezioni verso i due context (\code{GENERALEXCEPT} /
  \code{PGFAULTEXCEPT}) configurati nella sua \code{support\_t}. \\
\code{p6} & pass-up-or-die \emph{senza} support struct (syscall) & Esegue
  una syscall con codice positivo (\code{SYSCALL(1,...)}) \emph{senza}
  avere una support struct: il Nucleus deve terminarlo immediatamente
  (ramo ``die'' di \code{passUpOrDie}). Se \code{p6} sopravvive fino al
  \code{print} successivo, il test segnala un errore. \\
\code{p7} & pass-up-or-die \emph{senza} support struct (program trap) &
  Stesso principio di \code{p6} ma con un accesso a un indirizzo non
  valido invece di una syscall: verifica lo stesso ramo ``die'' per il tipo
  \code{GENERALEXCEPT} generato da un vero program trap. \\
\code{p8root} + \code{child1}/\code{child2} + 4 foglie & sotto-alberi
  profondi & Costruisce un albero di processi a 3 livelli (radice, 2 figli,
  4 nipoti) e lo ripete \textbf{4 volte di seguito}: verifica non solo che
  \code{terminateProcess} termini correttamente un sotto-albero intero
  (radice inclusa) quando la radice muore, ma anche che PCB e semafori
  usati vengano \emph{davvero} restituiti alle rispettive free list (la
  ripetizione fallirebbe rapidamente se ci fosse una perdita di risorse). \\
\code{p9}/\code{p10} & terminazione di un processo diverso da sé & \code{p9}
  crea \code{p10} e si blocca; \code{p10} termina \emph{il proprio padre}
  esplicitamente per PID (\code{TERMPROCESS} con \code{a1 = ppid}),
  verificando il ramo di \code{TERMPROCESS} in cui il chiamante non
  coincide col target (Sezione~\ref{subsub:sys2-termprocess} di questo
  documento: si veda \code{terminatingSelf}). \\
\code{hp\_p1}/\code{hp\_p2} & \code{YIELD} & Due processi ad alta priorità
  che cedono ripetutamente la CPU con \code{SYSCALL(YIELD,...)}: esercitano
  il meccanismo \code{yieldedProcess} dello scheduler
  (Sezione~\ref{sub:yield-scheduler}). \\
\bottomrule
\end{longtable}
\end{center}

Ogni sotto-processo stampa messaggi di progresso sul terminale 0 tramite una
funzione \code{print} interna (analoga a \code{termprint} di
\code{p1test.c}, ma qui implementata sopra una vera \code{DOIO}, dato che il
Nucleus è già disponibile), protetta da un semaforo di mutua esclusione
(\code{sem\_term\_mut}) per evitare che più processi scrivano
contemporaneamente e mescolino i caratteri sullo stesso terminale --- un
piccolo promemoria pratico che anche una risorsa condivisa ``interna'' come
il terminale di debug va protetta con un semaforo.

\paragraph{\code{p2test\_OLD.c}.} Nella cartella \code{phase2/} è presente
anche un file \code{p2test\_OLD.c}: è una versione precedente dello stesso
test driver, \textbf{non referenziata da nessun target CMake} (si veda
\code{CMakeLists.txt}: solo \code{phase2/p2test.c} compare nell'elenco
sorgenti del target \code{phase2}) e quindi non più compilata né eseguita;
resta nel repository solo come riferimento storico.

\end{document}
